
\textbf{Proposition 2.2.} \uline{Le complexe} @@M0@@ \uline{ne dépend pas}, \uline{à un isomorphisme canonique près dans la catégorie dérivée} @@M1@@, \uline{de la factorisation choisie par un} @@M2@@-\uline{schéma lisse}.

Soient @@M3@@ deux @@M4@@-schémas lisses par lesquels se factorise @@M5@@. Posons : @@M6@@; alors @@M7@@ est lisse sur @@M8@@, et aussi sur @@M9@@ et @@M10@@. De plus, les immersions @@M11@@ et @@M12@@ définissent une immersion @@M13@@ qui factorise @@M14@@. On peut donc se ramener à comparer les complexes @@M15@@ relatifs à deux factorisations par des @@M16@@-schémas lisses @@M17@@ et @@M18@@, tels qu'il existe un diagramme commutatif au-dessus de @@M19@@ :
@@M20@@
où @@M21@@ et @@M22@@ sont des immersions fermées, et @@M23@@ un morphisme lisse.

Par fonctorialité, on obtient alors des homomorphismes :
@@M24@@
d'où un diagramme :
@@M25@@

Comme @@M26@@ et @@M27@@ sont lisses sur @@M28@@, on obtient la suite exacte :
@@M29@@
Comme les Modules de différentielles considérés sont localement libres sur @@M30@@, on trouve en appliquant le foncteur @@M31@@ la suite exacte :
@@M32@@
